% !TeX root = ../main.tex
\section{Задача Коши и единственность}

\subsection*{Используемая теорема}
Для уравнения $y^{(n)}=f(x,y,y',\ldots,y^{(n-1)})$ непрерывность $f$
и её частных производных по фазовым переменным в окрестности
начальной точки обеспечивает существование и единственность
\emph{локального} решения для заданных $n$ начальных значений.
Это достаточное условие; его нарушение ещё не доказывает неединственность.
Графики в плоскости $(x,y)$ и траектории в расширенном пространстве
$(x,y,y',\ldots,y^{(n-1)})$ --- разные объекты.

\task{Ф. 225, а, г}{Области применимости теоремы}
\sourcepage{27}
\textbf{а)} $y'=2xy+y^2$. Здесь
\[
 f(x,y)=2xy+y^2,\qquad f_y=2x+2y
\]
непрерывны во всей $\R^2$. Через любую точку $(x_0,y_0)$ проходит
единственное локальное решение.

\textbf{г)} $y'=1+\tan y$. Здесь
\[
 f_y=\frac1{\cos^2y}.
\]
Обе функции непрерывны в полосах
$-\pi/2+k\pi<y<\pi/2+k\pi$, $k\in\Z$.
\answer{В пункте а подходят любые начальные данные; в пункте г требуется
$\cos y_0\ne0$. На прямых $y_0=\pi/2+k\pi$ само уравнение не определено.}
\editnote{В рукописи правая часть $1+\tan y$ ошибочно названа
непрерывной на всей плоскости.}

\task{Ф. 228, в, г}{Уравнения третьего порядка}
\sourcepage{28--29}
Обозначим начальные данные через
$y(x_0)=y_0$, $y'(x_0)=p_0$, $y''(x_0)=q_0$.

\textbf{в)} $(x-y)y'y'''=\ln(xy)$. Нормальная форма:
\[
 y'''=f(x,y,p,q)=\frac{\ln(xy)}{(x-y)p}.
\]
Она гладка в области $xy>0$, $x\ne y$, $p\ne0$.
Для полноты вычислим производные:
\[
 f_y=\frac1{y(x-y)p}+\frac{\ln(xy)}{(x-y)^2p},\qquad
 f_p=-\frac{\ln(xy)}{(x-y)p^2},\qquad f_q=0.
\]
\answer{Теорема гарантирует единственность при
$x_0y_0>0$, $x_0\ne y_0$, $p_0\ne0$; $q_0$ произвольно.}

\textbf{г)} $y''-yy'''=\sqrt[5]{y'-x}$. Нормальная форма:
\[
 f(x,y,p,q)=\frac{q-\sqrt[5]{p-x}}y.
\]
Она непрерывна при $y\ne0$. Для условия единственности нужны
\[
 f_y=-\frac{q-\sqrt[5]{p-x}}{y^2},\qquad
 f_p=-\frac1{5y\,|p-x|^{4/5}},\qquad f_q=\frac1y.
\]
\answer{Достаточные условия: $y_0\ne0$, $p_0\ne x_0$;
$q_0$ произвольно. При $y_0\ne0$, $p_0=x_0$ непрерывность всё ещё
даёт локальное существование, но приведённый критерий единственности не применим.}
Нельзя называть $y'=0$ условием «$y=\mathrm{const}$»:
нулевая производная в одной точке не означает постоянства функции.

\task{Ф. 229}{Могут ли пересекаться интегральные кривые?}
\sourcepage{29--30}
\textbf{а)} Для $y'=x+y^2$ функция $f=x+y^2$ и производная $f_y=2y$
непрерывны всюду. Два различных решения не могут пересекаться в общей
точке: одинаковые $x_0,y_0$ определяют одно локальное решение,
а единственность распространяет совпадение на общий интервал.

\textbf{б)} Для $y''=x+y^2$ единственность требует уже трёх чисел
$(x_0,y_0,y'_0)$. Можно задать, например,
\[
 y_1(x_0)=y_2(x_0)=y_0,\qquad y'_1(x_0)=1,\quad y'_2(x_0)=2.
\]
Оба локальных решения существуют, различны и пересекаются в $(x_0,y_0)$.
\answer{В пункте а --- нет; в пункте б --- да, с различными наклонами.}

\task{Ф. 230}{Могут ли различные решения касаться?}
\sourcepage{30}
Рассматриваются уравнения $y^{(n)}=x+y^2$ при $n=1,2,3$.
Касание графиков в точке означает совпадение $y$ и $y'$.

При $n=1$ уже совпадение $y$ при одном $x$ влечёт совпадение решений.
При $n=2$ совпавшие $y$ и $y'$ составляют полный набор данных Коши,
поэтому различные решения касаться не могут.

При $n=3$ зададим
\[
 y_1(x_0)=y_2(x_0)=y_0,\quad y'_1(x_0)=y'_2(x_0)=p_0,
 \quad y''_1(x_0)=1,\quad y''_2(x_0)=2.
\]
Теорема даёт два различных решения с общей касательной.
Различие видно также из разложения
$y_2(x)-y_1(x)=\tfrac12(x-x_0)^2+o((x-x_0)^2)$.
\answer{Для порядков 1 и 2 --- нет; для порядка 3 --- да.}

\task{Ф. 231}{Число решений при неполных или избыточных данных}
\sourcepage{31}
Для $y^{(n)}=x+y^2$ заданы $y(0)=1$, $y'(0)=2$.
\begin{itemize}
\item $n=1$: само уравнение требует $y'(0)=0+1^2=1$.
Условие $y'(0)=2$ противоречиво; решений нет.
\item $n=2$: заданы ровно два начальных значения. Существует
единственное локальное решение.
\item $n=3$: начальное значение $y''(0)=a$ можно выбрать произвольно.
Для каждого $a\in\R$ существует единственное локальное решение
с данными $(1,2,a)$; разные $a$ дают разные решения.
\end{itemize}
\answer{Соответственно: ноль, одно и бесконечно много локальных решений.}

\task{Ф. 233}{Две заданные функции как решения одного уравнения}
\sourcepage{31--32}
Пусть $y_1=x$, $y_2=x+x^4$. Требуется выяснить, при каком порядке
они могут быть решениями одного уравнения в нормальной форме,
удовлетворяющего условиям теоремы единственности.
В точке $x=0$:
\[
\begin{array}{c|ccccc}
 &y&y'&y''&y'''&y^{(4)}\\\hline
 y_1&0&1&0&0&0\\
 y_2&0&1&0&0&24
\end{array}
\]
Для $n\le4$ совпадают все $n$ начальных значений до производной
порядка $n-1$ включительно. Единственность тогда заставила бы
$y_1$ и $y_2$ совпасть, что неверно.
Для $n\ge5$ обе функции действительно удовлетворяют одному и тому же
гладкому уравнению $y^{(n)}=0$.
\answer{Возможно в точности при $n\ge5$.}
\editnote{Отсутствие противоречия с теоремой ещё не доказывает
существование нужного уравнения. Поэтому добавлен явный пример
$y^{(n)}=0$. Без предположения об условиях единственности вывод для
малых порядков в такой форме делать нельзя.}

\task{Д. 3, а}{Непродолжаемое решение}
\sourcepage{32}
\[
 y'=-y^2,\qquad y(1)=-1.
\]
Нулевое решение не подходит под начальные данные.
Для ненулевого решения
\[
 \frac{\dd y}{y^2}=-\dd x,\qquad \frac1y=x+C,
 \qquad -1=1+C.
\]
\answer{$\displaystyle y(x)=\frac1{x-2}$ на $(-\infty,2)$.}
Хотя выражение определено также на $(2,+\infty)$, этот второй интервал
не содержит начальной точки. Решение задачи Коши задаётся на связном
интервале; писать его область как $\R\setminus\{2\}$ некорректно.
При $x\to2-0$ функция неограниченна, поэтому продолжения нет.

Все интегральные кривые уравнения: $y=0$ и $y=1/(x-C)$ на интервалах
по одну сторону от $C$. Поскольку $f_y=-2y$ непрерывна всюду,
через каждую точку проходит единственное локальное решение.
Особого решения в смысле нарушения единственности или огибающей здесь нет;
$y=0$ является отдельным равновесным решением, потерянным при делении.

\begin{center}
\begin{tikzpicture}[x=1cm,y=.65cm]
 \draw[->] (-2.5,0)--(4.1,0) node[right] {$x$};
 \draw[->] (0,-3.2)--(0,3.2) node[above] {$y$};
 \draw[dashed,gray] (2,-3)--(2,3);
 \draw[thick,blue!65!black,domain=-2.4:1.65,samples=80] plot(\x,{1/(\x-2)});
 \draw[gray,domain=2.35:3.9,samples=60] plot(\x,{1/(\x-2)});
 \fill (1,-1) circle(2pt) node[below left] {$(1,-1)$};
 \node[above right] at (2,0) {$2$};
\end{tikzpicture}
\end{center}

\task{Д. 3, б}{Решение с нулевым участком}
\sourcepage{32--33}
\[
 y'=3\sqrt[3]{y^2},\qquad y(-4)=-1,\quad y(2)=1.
\]
На ненулевых ветвях $y=(x-C)^3$. Левая начальная точка определяет
$C=-3$, правая --- $C=1$. Отрицательная ветвь достигает нуля при
$x=-3$, положительная начинается при $x=1$.
\answer{Единственное решение с обоими указанными значениями:
\[
 y(x)=\begin{cases}
 (x+3)^3,&x<-3,\\
 0,&-3\le x\le1,\\
 (x-1)^3,&x>1,
 \end{cases}\qquad x\in\R.
\]}
Обоснование полноты такое же, как в задаче № 55: $y'\ge0$, поэтому
между первым и последним нулями функция обязана быть нулевой.
Вне них константы кубических ветвей уже определены условиями.
Производная $f_y=2/\sqrt[3]y$ существует и непрерывна лишь при $y\ne0$.
Поэтому одна начальная точка вне оси задаёт единственную локальную
ветвь, но после достижения нуля возможно произвольное время ожидания.
Второе условие фиксирует это время. Решение $y=0$ особое: в каждой его
точке имеются другие решения с тем же начальным значением.

\begin{center}
\begin{tikzpicture}[x=.9cm,y=.65cm]
 \draw[->] (-4.6,0)--(2.7,0) node[right] {$x$};
 \draw[->] (0,-2.7)--(0,2.7) node[above] {$y$};
 \draw[thick,blue!65!black,domain=-4.3:-3,samples=50] plot(\x,{(\x+3)^3});
 \draw[very thick,blue!65!black] (-3,0)--(1,0);
 \draw[thick,blue!65!black,domain=1:2.3,samples=50] plot(\x,{(\x-1)^3});
 \fill (-4,-1) circle(2pt) node[below left] {$(-4,-1)$};
 \fill (2,1) circle(2pt) node[above right] {$(2,1)$};
 \node[below] at (-3,0) {$-3$}; \node[below] at (1,0) {$1$};
\end{tikzpicture}
\end{center}
